
\begin{center}
    \includegraphics[width=9.5cm]{../../assets/logo7ZPUs.jpg}\\
    \small\hspace{10cm} 7zpus.swe@gmail.com\\
    \vspace{0.5cm}
    \LARGE \textbf{Specifica Tecnica}\\
    \vspace{0.3cm}
\end{center}

\vspace{0.3cm}
\hrule
\vspace{0.5cm}


\section*{Tabella di Versionamento}
\begin{table}[H]
    \begin{adjustwidth}{-4cm}{-4cm}
    \centering
    \begin{tabular}{|c|c|c|c|c|}
        \hline
        \textbf{Versione} & \textbf{Data} & \textbf{Autore}  & \textbf{Verificatore} & \textbf{Descrizione} \\
        \hline
        0.2.0 & 2026/04/08 & Georgescu Diana & Soligo Lorenzo & Aggiunta descrizione modulo Export \\
        \hline
        0.1.0 & 2026/02/28 & Laoud Zakaria & Georgescu Diana & \begin{tabular}[c]{@{}c@{}} Creazione del documento\\ e stesura struttura iniziale \end{tabular} \\
        \hline
    \end{tabular}
    \end{adjustwidth}
\end{table}

\newpage
\tableofcontents

\newpage
\listoffigures

\listoftables

\newpage
\section{Introduzione}

\subsection{Obiettivo del documento}
L’obiettivo del presente documento è descrivere in modo rigoroso le scelte tecnologiche e l’architettura del prodotto \textit{DIPReader}, adottando il modello di progettazione \textit{C4}.\\
L’analisi verrà condotta partendo dalla descrizione del contesto applicativo, per poi approfondire la struttura dei container, delle componenti e, infine, del codice, al fine di fornire una rappresentazione completa, 
coerente e sistematica del sistema.

\subsection{Glossario}
Ogni termine tecnico o con un significato particolare nell’ambito dell’Ingegneria del Software, utilizzato nella documentazione di progetto, è definito nell’apposito documento
\href{https://7-zpus.github.io/Docs/assets/3_PB/DocumentiInterni/Glossario.pdf}{Glossario }\ped{(v2.0.0)}.

\subsection{Riferimenti}

    \subsubsection{Riferimenti normativi}
    \begin{itemize}
        \item \href{https://github.com/7-ZPUs/Docs/blob/main/3_PB/DocumentiInterni/NormeDiProgetto.pdf}{Norme di Progetto v2.0.0 }\ped{(v2.0.0)}
        \item \href{https://en.wikipedia.org/wiki/Software_design_description}{Standard SDD/IEEE/1016:2009}\ped{(ultimo accesso: 02/03/2026)}
        \item \href{https://www.math.unipd.it/~tullio/IS-1/2025/Dispense/PD1.pdf}{Regolamento di Progetto Didattico a.a.
            2025/2026} \ped{(ultimo accesso: 02/03/2026)}
    \end{itemize}

    \subsubsection{Riferimenti informativi}
    \begin{itemize}
        \item \href{https://www.math.unipd.it/~tullio/IS-1/2025/Progetto/C3.pdf}{Capitolato C3: DIPReader} \ped{(ultimo accesso: 10/02/2026)}
        \item \href{https://www.math.unipd.it/~rcardin/swea/2022/Design%20Pattern%20Architetturali%20-%20Dependency%20Injection.pdf}
        {Dispense su Dependency Injection }\ped{(Ultimo accesso: 02/03/2026)}
        \item \href{https://www.math.unipd.it/~rcardin/swea/2023/Object-Oriented%20Progamming%20Principles%20Revised.pdf}
        {Dispense su OOP }\ped{(Ultimo accesso: 02/03/2026)}
        \item \href{https://www.math.unipd.it/~rcardin/swea/2023/Diagrammi%20delle%20Classi.pdf}
        {Dispense su Diagrammi delle classi }\ped{(Ultimo accesso: 02/03/2026)}
        \item \href{https://www.math.unipd.it/~rcardin/swea/2022/Software%20Architecture%20Patterns.pdf}
        {Dispense su Pattern architetturali }\ped{(Ultimo accesso: 02/03/2026)}
        \item \href{https://www.math.unipd.it/~rcardin/swea/2022/Design%20Pattern%20Creazionali.pdf}
        {Dispense su Pattern creazionali }\ped{(Ultimo accesso: 02/03/2026)}
        \item \href{https://www.math.unipd.it/~rcardin/swea/2022/Design%20Pattern%20Strutturali.pdf}
        {Dispense su Pattern strutturali }\ped{(Ultimo accesso: 02/03/2026)}
        \item \href{https://www.math.unipd.it/~rcardin/swea/2021/SOLID%20Principles%20of%20Object-Oriented%20Design_4x4.pdf}
        {Dispense su Principi SOLID }\ped{(Ultimo accesso: 02/03/2026)}
        \item \href{https://www.math.unipd.it/~tullio/IS-1/2023/Dispense/T6.pdf}
        {Dispense sulla Progettazione software }\ped{(Ultimo accesso: 02/03/2026)}
        \item \href{https://www.math.unipd.it/~tullio/IS-1/2023/Dispense/T7.pdf}
        {Dispense sulla Qualità del software }\ped{(Ultimo accesso: 02/03/2026)}
        \item \href{https://github.com/rcardin/hexagonal}
        {Repository Architettura esagonale (Kotlin) }\ped{(Ultimo accesso: 02/03/2026)}
        \item \href{https://github.com/rcardin/hexagonal-java/}
        {Repository Architettura esagonale (Java) }\ped{(Ultimo accesso: 02/03/2026)}
        \item \href{https://github.com/rcardin/swe-imdb}
        {Repository Ingegneria del software - swe-imdb }\ped{(Ultimo accesso: 02/03/2026)}
        \item \href{https://c4model.com/}
        {C4 Model }\ped{(Ultimo accesso: 03/03/2026)}
        \item \href{https://plantuml.com/}
        {PlantUML }\ped{(Ultimo accesso: 03/03/2026)}
        \item Riferimenti a scelte tecnologiche:
        \begin{itemize}
            \item \href{https://www.typescriptlang.org/}
            {TypeScript }\ped{(Ultimo accesso: 03/03/2026)}
            \item \href{https://angular.io/}
            {Angular }\ped{(Ultimo accesso: 03/03/2026)}
            \item \href{https://www.electronjs.org/}
            {Electron }\ped{(Ultimo accesso: 03/03/2026)}
            \item \href{https://www.sqlite.org/}
            {SQLite }\ped{(Ultimo accesso: 03/03/2026)}
            \item \href{https://github.com/asg017/sqlite-vss}
            {sqlite-vss }\ped{(Ultimo accesso: 03/03/2026)}
            \item \href{https://onnxruntime.ai/}
            {ONNX Runtime }\ped{(Ultimo accesso: 03/03/2026)}
            \item \href{https://huggingface.co/docs/transformers.js}
            {Transformers.js }\ped{(Ultimo accesso: 03/03/2026)}
            \item \href{https://nodejs.org/}
            {Node.js }\ped{(Ultimo accesso: 03/03/2026)}
        \end{itemize}
        \item \href{https://7-zpus.github.io/Docs/assets/3_PB/DocumentiInterni/Glossario.pdf}{Glossario }\ped{(v2.0.0)}
        \item \href{https://github.com/7-ZPUs/Docs/blob/main/3_PB/DocumentiEsterni/PianoDiProgetto.pdf}{Piano di Progetto }\ped{(v2.0.0)}
        \item \href{https://github.com/7-ZPUs/Docs/blob/main/3_PB/DocumentiEsterni/PianoDiQualifica.pdf}{Piano di Qualifica }\ped{(v2.0.0)}
        \item Verbali interni:
        \item Verbali esterni:
    \end{itemize}
    
    

\newpage
\section{Tecnologie utilizzate}
\newpage
\section{Contesto del sistema}

\subsection{Modello architetturale}

\subsection{Descrizione degli attori esterni}

\subsection{Descrizione dei sistemi esterni}
